% Source: https://artofproblemsolving.com/community/c3h2241851_propability
% Stable identity: aops-topic-2241851-q01
\textbf{(a)} A split is determined by the set of $4$ children in group $A$, and there are $\dbinom84=70$ equally likely choices. The boys are together exactly when group $A$ consists of the $4$ boys or of the $4$ girls: $2$ choices. So the probability is
\[\dfrac{2}{70}=\boxed{\dfrac{1}{35}}.\]
(Check: whatever group the first boy is in, the other three boys must take the $3$ remaining places of that group among the $7$ remaining places: $\dfrac{3}{7}\cdot\dfrac{2}{6}\cdot\dfrac{1}{5}=\dfrac{1}{35}$.)

\textbf{(b)} The $16$ patterns are not equally likely, because the boys are not placed independently: each group has exactly $4$ places, so every boy placed in $A$ leaves fewer places there for the others. A pattern with $k$ boys in $A$ happens only if the remaining $4-k$ places of $A$ go to girls, which can happen in $\dbinom{4}{4-k}$ of the $70$ splits. So $AAAA$ has probability $\dfrac{1}{70}$, not $\dfrac{1}{16}$, while a pattern such as $AABB$ has probability $\dfrac{6}{70}$. (The student's model would be right only if each child chose a group by a coin toss, with no limit on the group sizes.)
